\scanpage{23}
\mspage{10}
Le sous-groupe $\mathcal E_\tau$ de $\mathcal E$ image réciproque de $\mathbb Z/2$ engendré par $\tau$, est le groupe cartographique ``triangulé \unsure{pondéré|pointé}'' de Malgoire, ayant comme gén $\tau_0,\tau_1,\tau_\infty$ \deleted{\unsure{relati.}} avec
\[
(17)\quad\left\{\begin{aligned}
\deleted{\tau_1\tau_\infty={}}\ l_0&=\tau_{\deleted{\unsure{0}}\infty}\tau_{\deleted{\unsure{\infty}}1}\\
\deleted{\tau_\infty\tau_0={}}\ l_1&=\tau_{\deleted{\unsure{\infty}}0}\tau_\infty\\
l_\infty&=\tau_1\tau_0
\end{aligned}\right.
\]
\margin{\uline{NB} $-$}
$\mathcal E_\tau$ peut être considéré comme le groupe ayant comme \inserted{gén. les $\tau_i$, avec les} seules relations (16) --- on a un hom signature évident $\tau_i\longrightarrow-1$ de $\mathcal E_\tau$ dans $\pm1$, dont le noyau est $\pi$ engendré par les $l_i$ décrits par (17)\ldots

\figure{\begin{tikzpicture}\draw (0,0)--(0.8,0.025);\end{tikzpicture}}
\margin{On a les relations $\unsure{\rho}$}
Revenons à une description plus symétrique des autom. extérieurs à lacets de $\mathcal E$. On désignera par $\hat{\mathcal E}$ le complété profini de $\mathcal E$, extension de $\hat\Gamma_{0,3}=\Gamma_{0,3}=\unsure{\mathfrak S_3|\mathcal S_3}\times\mathbb Z/2$ par $\hat\pi$, et par $\widehat{\widehat{\mathcal E}}$ le groupe des autom. \deleted{extérieurs} à lacets de $\hat\pi$, extension de $\widehat{\widehat\Gamma}=\operatorname{Aut}\operatorname{ext}_{\mathrm{lac}}(\hat\pi)$ par $\hat\pi$
\[
(18)\qquad
\diagram{\begin{tikzcd}
1\arrow[r]&\hat\pi\arrow[r]&\widehat{\widehat{\mathcal E}}\arrow[r]&\widehat{\widehat\Gamma}\arrow[r]\arrow[d,equal]&1\\
&&&\operatorname{Aut}\operatorname{ext}_{\mathrm{lac}}(\hat\pi)&
\end{tikzcd}}
\]
\mspage{22}
